\documentclass{article}
\usepackage{titlesec}
\usepackage[utf8]{inputenc}
\usepackage[a4paper, total={6in, 9in}]{geometry}
\usepackage{fancyhdr}
\usepackage{graphicx} % Add the graphicx package
\usepackage{lipsum} % For placeholder text
\usepackage{tikz}
\usepackage{cite}
\usepackage{bm}
\usepackage{acronym}
\usepackage{amsmath,amssymb,amsfonts}
\usepackage{algorithmic}
\usepackage{graphicx}
\usepackage{textcomp}
\usepackage{xcolor}
\usepackage[numbers]{natbib}
\bibliographystyle{ieeetrann}

% Set up page headers and footers
\pagestyle{fancy}
\fancyhf{}
\rhead{Group No}
\lhead{Project Title}
\cfoot{\thepage}

% Redefine the abstract environment to be single column
\renewenvironment{abstract}
{\par\noindent\textbf{\abstractname}\ \ignorespaces}{\par}

% Redefine section format
\titleformat{\section}
{\normalfont\Large\bfseries}{\thesection}{1em}{}

\begin{document}
	
	% Cover Page
	\begin{titlepage}
		
		\centering
		\vspace*{0.5cm}
		\includegraphics[width=3cm]{logo.png} % Replace with the path to your university logo
		\par\vspace{0.02cm}
		Department of Electronic \& Telecommunication Engineering,
		University of Moratuwa, Sri Lanka.
		\par\vspace{2cm}
		{\LARGE\bfseries Analog Computer\par}
		\vspace{7cm}
		{\large Group Members:\par}
		\begin{tabular}{c c}
			& \\
	 240445H	&	N.Naagul\\
Member 2 Index No	&	Member 2 Name\\
Member 3 Index No	&	Member 3 Name\\
Member 4 Index No	&	Member 4 Name \\
		\end{tabular}\\
		\vspace{1.5cm}
		{Submitted in partial fulfillment of the requirements for the module\par}
		{EN 2091 Laboratory Practice and Projects\par}
	
		\vspace{1.0cm}
		{\large {\today}\par}
		\vfill
	\end{titlepage}
	
	\newpage
	
	% Abstract
	\begin{abstract}
		
	\end{abstract}
	
	

	
	\newpage
	
\section*{Abbreviations and Acronyms}

\begin{acronym}
	\acro{API}{Application Programming Interface}
	\acro{HTML}{Hypertext Markup Language}
	\acro{CSS}{Cascading Style Sheets}
	\acro{PDF}{Portable Document Format}
\end{acronym}
	\newpage
	% Table of Contents
	\tableofcontents
	\newpage
	
	% Sections
	\section{Introduction and Functionality}

	
	
	\section{System Architecture}
This is a sample LaTeX document that demonstrates how to include figures, tables, and references.



\begin{figure}[htbp]
	\centering
	\includegraphics[width=0.6\textwidth]{example-image-a}
	\caption{Example Figure.}
	\label{fig:example}
\end{figure}
\subsection{Table and Figure}
In Figure \ref{fig:example}, we can see an example figure. Table \ref{tab:example} shows an example table.\\
\begin{table}[htbp]
	\centering
	\caption{Example Table.}
	\begin{tabular}{|l|l|}
		\hline
		\textbf{Header 1} & \textbf{Header 2} \\
		\hline
		Row 1, Col 1 & Row 1, Col 2 \\ \hline
		Row 2, Col 1 & Row 2, Col 2 \\ \hline
			\end{tabular}
		\label{tab:example}
\end{table} 
\subsection{Equations}
Number your equations consecutively and punctuate equations with commas or periods when they are integrated into a sentence. The system model is given in equation~\eqref{eq}
\begin{equation}
	a + b = \gamma \label{eq}.
\end{equation}
An example of multi-line equation is given below 
\begin{equation}
\begin{aligned}
	\bm{a} &= \bm{b} + \bm{c}, \\
	&= \left(\bm{1}+\bm{\gamma}\right)^2 + \bm{\delta}^3.
\end{aligned}
\label{eq1} 
\end{equation}
For vectors, represent them in bold using lowercase letters, such as \textbf{v}. For matrices, use bold uppercase letters, like \textbf{M}. Further, when representing vectors and matrices denote their dimensions explicitly, such as $\mathbf{v} \in \mathbb{R}^{N \times 1}$, where $\mathbb{R}$ represents the real numbers. Similarly, for matrices, express the dimensions, e.g., $\mathbf{M} \in \mathbb{R}^{m \times n}$, indicating $m$ rows and $n$ columns. Define abbreviations and acronyms the first time they are used in the text.
\subsection{Guidelines for References}
Include the reference list in the "References" section at the end of the document.  When referring to a source, use the reference number as shown below.\\
The invention of the transistor in 1947 by Bardeen, Brattain, and Shockley revolutionized electronic technology, enabling the development of smaller, more efficient electronic devices~\cite{bardeen1948transistor}. In electronic circuits, operational amplifiers (op-amps) play a pivotal role due to their versatile applications, such as signal amplification and filtering \cite{sedra_smith_microelectronic_circuits}.
	
	

	\section{Component Selection}
	
	
	\section{PCB Design}
	
	
	\section{Enclosure Design}
	
	\section{Software Simulation and Hardware Testing}

	
	\section{Conclusion \& Future Works}
	
	\section{Contribution of Group Members}
	
	
	\section*{Acknowledgment}

	
	
\bibliography{ref}

\bibliographystyle{IEEEbib}
	
	\section*{Appendix}

	
\end{document}
